\begin{question}{19}{
    De Koninklijke Marine onderzoekt op haar luchtverdedigings- en commandofregatten (LCF) de kwetsbaarheden van het ECDIS-navigatiesysteem op het gebied van GPS-spoofing.
    Tijdens een oefening op de Noordzee wordt de gemiddelde afwijking tussen de werkelijke positie en de door ECDIS weergegeven positie gemeten (in meters).

    Volgens de technische handboeken zou de gemiddelde afwijking onder normale omstandigheden minder dan $12$ meter moeten zijn.
    Bij een steekproef van achttien metingen wordt een gemiddelde van $14.1$ meter en een standaardafwijking van $1.4$ meter waargenomen.
    Neem aan dat de afwijkingen normaal verdeeld zijn.
}
    \subquestion{6}{
        Voer een geschikte hypothesetoets uit om te onderzoeken of de gemiddelde afwijking voldoet aan de claim in de technische handboeken.
    }
    \solution{

    }
    
    \subquestion{6}{
        Aangezien precieze lokalisering van essentieel operationeel belang is om de veiligheid van de vloot te waarborgen, wil men de gemiddelde afwijking in posities op 1 meter nauwkeurig kunnen schatten.
        Wat is de minimale steekproefomvang $n$ waarbij het \SI{95}{\percent}-betrouwbaarheidsinterval voor $\mu$ een afwijking heeft van $\pm 1$ meter.
    }
    \solution{
        
    }

    

\end{question}